\section{Conclusiones}

Las cargas críticas de la vivienda, refrigerador, congelador, iluminación, cargadores y módem, consumen \SI{4.33}{kWh} por día y piden \SI{783}{W} si todas funcionan a la vez. Con el método del curso, dos días de autonomía, \SI{24}{V} y 70\,\% de descarga, el banco requiere \SI{758}{Ah}, que se cubren con cuatro baterías Power Queen LiFePO\textsubscript{4} de \SI{200}{Ah} en paralelo, el máximo que admite el fabricante.

La cara este del techo, de \SI{12}{\degree}, recibe \SI{1720}{kWh/m^2} al año, 6.5\,\% más que la cara oeste; el plano óptimo recibiría 1.7\,\% más. Con los \SI{3.73}{HSP} de noviembre, el peor mes, se requieren cuatro módulos AIKO de \SI{450}{W}, que se conectan en dos cadenas de dos para cumplir los límites de tensión del regulador entre 7 y \SI{63.2}{\celsius}.

El regulador Victron SmartSolar 150/85 cumple la corriente de \SI{75}{A} que pide el método del curso, y el inversor-cargador Victron MultiPlus-II de \SI{3000}{VA} trabaja al 33\,\% de su potencia continua y al 32\,\% de su pico; si falla la red, conmuta las cargas a su inversor en menos de \SI{20}{ms}.

La energía disponible para las cargas supera la demanda todos los meses: en noviembre la cobertura es de 116\,\%, y en el año hay \SI{2319}{kWh} disponibles para \SI{1582}{kWh} de demanda.

El sistema cuesta unos ¢7.10 millones, de los cuales el 48\,\% son las baterías. Reduce la factura en ¢128\,165 al año (54\,\%), pero con el mantenimiento y los reemplazos la inversión no se recupera en 25 años (VAN de ¢8.00 millones negativos). Su justificación es la continuidad del servicio de las cargas críticas, no el ahorro.

Antes de construir deben confirmarse los precios con distribuidores locales, las distancias de los tramos de cableado, la separación y el estado de las clavadoras del techo y la reducción de potencia del regulador por la altitud del sitio.
